% ============================================================================
%  Directional signal: daily modulation from crystal anisotropy
%  Figure file: directional_modulation.png (NuFog_v2/outputs)
%  Bibliography entries are at the end of this file (move them to your .bib).
% ============================================================================

\section{Directional signal: daily modulation from crystal anisotropy}
\label{sec:directionality}

The rates presented so far were computed for an isotropic response, obtained by
averaging the phonon density of states over directions. Most of the crystals
considered here are not cubic, however, and their lattice response depends on the
direction of the momentum transfer $\vec q$. Because the dark-matter (DM) velocity
distribution seen in the laboratory is itself anisotropic --- the DM ``wind'' arrives
from the direction opposite to the Earth's velocity $\vec v_E$, which points towards the
constellation Cygnus --- the interaction rate depends on the orientation of the crystal
with respect to $\vec v_E$. The Earth's rotation changes this orientation with a period
of one sidereal day, producing a daily modulation of the rate. Neutrinos do not share
this signature: solar neutrinos follow the Sun and modulate with the solar day, while the
atmospheric and diffuse supernova components are approximately isotropic. A daily
modulation locked to sidereal time is therefore a property of the DM signal alone.

\subsection{Anisotropic response in the incoherent approximation}

In the incoherent approximation the direction dependence enters only through the
per-atom phonon density-of-states tensor $D_d^{ij}(\omega)$, obtained from the phonon
eigenvectors $\vec e_{d,\nu}(\vec k)$ as
\begin{equation}
  D_d^{ij}(\omega) = \frac{1}{N_k}\sum_{\nu,\vec k}
  e^{i\,*}_{d,\nu}(\vec k)\, e^{j}_{d,\nu}(\vec k)\;
  \delta\!\left(\omega-\omega_{\nu}(\vec k)\right),
  \qquad
  \int d\omega\, D_d^{ij}(\omega) = \delta^{ij} ,
\end{equation}
so that $\frac13\mathrm{Tr}\,D_d(\omega)$ is the unit-normalized scalar density of states
used in Sec.~\ref{sec:multiphonon}. For a momentum transfer along $\hat q$, the relevant
quantity is the projected density of states
$D_d^{\hat q}(\omega) = \hat q_i D_d^{ij}(\omega)\hat q_j$, which replaces the isotropic
$D_d(\omega)$ everywhere in the multiphonon expansion:
\begin{equation}
  W_d(\vec q) = \frac{q^2}{4 m_d}\int d\omega\,\frac{D_d^{\hat q}(\omega)}{\omega},
  \qquad
  T^{\hat q}_{d,1}(\omega) = \frac{D_d^{\hat q}(\omega)}{\omega},
  \qquad
  T^{\hat q}_{d,n} = T^{\hat q}_{d,1} * T^{\hat q}_{d,n-1} .
\end{equation}
Since $D_d^{\hat q}$ is quadratic in $\hat q$, the structure factor satisfies
$S(\vec q,\omega) = S(-\vec q,\omega)$. For the uniaxial (tetragonal, trigonal) crystals
studied here it depends, to a good approximation, only on the polar angle $\theta_q$
between $\hat q$ and the crystallographic $c$ axis. We verified this directly from the
tensors: the azimuthal variation of the species-summed $\langle 1/\omega\rangle$ is
zero for CaWO$_4$ and LiNbO$_3$ and 2--6\% for Al$_2$O$_3$ and TeO$_2$, compared with
polar anisotropies of up to 48\% (W in CaWO$_4$). The remaining azimuthal dependence is
averaged out. Silicon is cubic, its tensor is isotropic, and it serves as a null test.

\subsection{Rate for an anisotropic velocity distribution}

For an isotropic response the velocity integral reduces to the mean inverse speed
$\eta(v_{\min})$. In general one needs the Radon transform of the velocity distribution,
\begin{equation}
  g(w,\hat q) = \int d^3v\, f(\vec v)\,\delta\!\left(\vec v\cdot\hat q - w\right),
\end{equation}
which for the truncated Maxwellian of the Standard Halo Model, boosted by $\vec v_E$, is
\begin{equation}
  g(w,\hat q) = \frac{1}{N_{\rm esc}\sqrt{\pi}\,v_0}
  \left[ e^{-(w+\hat q\cdot\vec v_E)^2/v_0^2} - e^{-v_{\rm esc}^2/v_0^2} \right]
  \Theta\!\left(v_{\rm esc}-|w+\hat q\cdot\vec v_E|\right),
  \quad
  N_{\rm esc} = \mathrm{erf}(z) - \frac{2z}{\sqrt\pi}e^{-z^2},
\end{equation}
with $z=v_{\rm esc}/v_0$. The differential rate becomes
\begin{equation}
  \frac{dR}{d\omega} = \frac{1}{\rho_T}\frac{\rho_\chi}{m_\chi}
  \frac{\sigma_n}{4\pi\mu_{\chi n}^2}
  \int q\,dq \int\frac{d\Omega_q}{4\pi}\;
  2\,g\!\left(v^*, \hat q\right) S(\vec q,\omega),
  \qquad
  v^* = \frac{q}{2m_\chi} + \frac{\omega}{q},
  \label{eq:dRdw_dir}
\end{equation}
which reduces to the isotropic expression because
$\int (d\Omega_q/4\pi)\, 2 g(w,\hat q) = \eta(w)$. Equation~\eqref{eq:dRdw_dir} depends
on the angle $\theta_E$ between $\vec v_E$ and the $c$ axis, with
$R(\theta_E) = R(\pi-\theta_E)$.

Numerically, $S$ is tabulated on 12 independent Gauss--Legendre nodes in
$\cos\theta_q$, using $S(\theta_q)=S(\pi-\theta_q)$. Because $g$ has a kink at the
escape velocity, its angular integral is computed by product quadrature: $S$ is held
fixed within each node's cell in $\cos\theta_q$, while $g$ is sampled finely and averaged
over the azimuth of $\hat q$. This reproduces $\eta(v_{\min})$ to better than
$10^{-3}$ wherever $\eta$ exceeds $10^{-3}$ of its maximum. The rate is evaluated for ten
values of $\theta_E \in [0^\circ,90^\circ]$ and binned exactly as in the isotropic
pipeline, so the same thresholds and energy resolution (Sec.~\ref{sec:resolution}) can
be applied.

Two consistency checks were performed. First, averaging the directional rate over all
orientations of $\vec v_E$ reproduces the isotropic rate. For Si the agreement is better
than $10^{-4}$ at all masses. For CaWO$_4$, the most anisotropic material, the residual
is at most $0.3\%$. This residual is physical: the average of $S(D^{\hat q})$ differs
from $S$ evaluated with the averaged tensor at second order in the anisotropy, because
the Debye--Waller factor and the multiphonon convolutions are nonlinear in $D$. It also
shows that the isotropic discovery curves of Sec.~\ref{sec:materials} describe the
time-averaged rate to within $0.3\%$. Second, the cubic Si crystal shows no modulation,
as required by symmetry.

\subsection{Laboratory geometry}

For a crystal axis with equatorial coordinates $(\alpha_c,\delta_c)$,
\begin{equation}
  \cos\theta_E = \sin\delta_c\sin\delta_E
  + \cos\delta_c\cos\delta_E\cos(\alpha_c-\alpha_E),
\end{equation}
where $(\alpha_E,\delta_E)\simeq(21^{\rm h}12^{\rm m},\,+48^\circ)$ is the direction of
$\vec v_E$. A laboratory-fixed axis rotates with the local sidereal time (LST). For the
$c$ axis mounted vertically at LNGS (latitude $42.45^\circ$\,N),
$\delta_c = 42.45^\circ$ and $\alpha_c = \mathrm{LST}$, so $\theta_E$ sweeps from
$6^\circ$ to $90^\circ$ every sidereal day (Fig.~\ref{fig:directional}a). This covers
almost the full independent range of $\theta_E$ and is close to the optimal
orientation. A horizontal $c$ axis pointing east only sweeps $48^\circ$--$90^\circ$.

\begin{figure}[t]
  \centering
  \includegraphics[width=\textwidth]{directional_modulation.png}
  \caption{Daily modulation of the DM rate from crystal anisotropy (heavy mediator,
  $\sigma=E_{\rm th}/5$ resolution, Standard Halo Model).
  (a) Angle between the crystal $c$ axis and the DM wind over a sidereal day at LNGS,
  for a vertical and a horizontal (east-pointing) $c$ axis.
  (b) Anisotropy of the rate above threshold,
  $[R(0^\circ)-R(90^\circ)]/[R(0^\circ)+R(90^\circ)]$, as a function of DM mass for
  $E_{\rm th}=30$\,meV (solid) and $100$\,meV (dashed).
  (c) Rate over a sidereal day, normalized to its mean, for $m_\chi = 1$\,MeV,
  $E_{\rm th}=100$\,meV, and a vertical $c$ axis.}
  \label{fig:directional}
\end{figure}

\subsection{Results}

The anisotropy of the rate is shown as a function of DM mass in
Fig.~\ref{fig:directional}b, and the corresponding daily modulation in
Fig.~\ref{fig:directional}c. Table~\ref{tab:directional} summarizes the values on the
low-mass plateau. Three features stand out.

\begin{enumerate}
  \item \textbf{The modulation is confined to $m_\chi \lesssim 50$\,MeV.}
  At larger masses the rate is dominated by momentum transfers in the impulse regime,
  where the response approaches that of a free nucleus and is isotropic. The modulation
  therefore probes the same mass range in which multiphonon emission controls the
  sensitivity.

  \item \textbf{Its size is at the percent level and depends on the material.}
  At $E_{\rm th}=100$\,meV, LiNbO$_3$ and TeO$_2$ show the largest anisotropy ($\simeq
  4\%$), followed by Al$_2$O$_3$ ($2.4\%$) and CaWO$_4$ ($1.2\%$). The ordering follows
  the anisotropy of the light atoms (O, Li), which dominate the multiphonon response at
  low mass. The large W anisotropy in CaWO$_4$ shows up mainly at intermediate masses,
  around $20$\,MeV, where the heavy nucleus contributes more of the rate.

  \item \textbf{The sign can reverse with threshold.}
  For LiNbO$_3$, Al$_2$O$_3$ and CaWO$_4$, the rate at $\theta_E=0^\circ$ is larger
  than at $90^\circ$ for $E_{\rm th}=30$\,meV and smaller for $E_{\rm th}=100$\,meV.
  Events near a low threshold come mostly from the low-frequency part of
  $D^{\hat q}(\omega)$, weighted by $1/\omega$, and favour the direction with softer
  modes. Events above the highest optical phonon require several phonons and draw on the
  high-frequency part of the spectrum, so they favour the direction with stiffer modes.
  A daily phase that reverses between energy windows is a distinctive signature that a
  background is unlikely to mimic.
\end{enumerate}

\begin{table}[t]
  \centering
  \caption{Anisotropy of the DM rate above threshold,
  $[R(0^\circ)-R(90^\circ)]/[R(0^\circ)+R(90^\circ)]$, at $m_\chi = 1$\,MeV, with
  $\sigma=E_{\rm th}/5$. For a vertical $c$ axis at LNGS the peak-to-peak daily
  variation is approximately twice this value.}
  \label{tab:directional}
  \begin{tabular}{lcc}
    \hline\hline
    Material & $E_{\rm th}=30$\,meV & $E_{\rm th}=100$\,meV \\
    \hline
    CaWO$_4$  & $+0.1\%$ & $-1.2\%$ \\
    Al$_2$O$_3$ & $+0.7\%$ & $-2.4\%$ \\
    TeO$_2$   & $+1.6\%$ & $+4.0\%$ \\
    LiNbO$_3$ & $+1.6\%$ & $-4.3\%$ \\
    Si        & $0$      & $0$      \\
    \hline\hline
  \end{tabular}
\end{table}

A sidereal modulation of fractional amplitude $A$ is detected at $3\sigma$ with roughly
$N \approx 18/A^2$ signal events. For $A = 4\%$ this is of order $10^4$ events, far more
than the $\mathcal{O}(10)$ events needed for a discovery above the neutrino background
(Sec.~\ref{sec:discovery}). Directionality is therefore a tool to \emph{confirm} the DM
origin of an excess --- through its sidereal period, its phase, and the threshold
dependence of that phase --- rather than to extend the discovery reach.

\subsection{Comparison with the literature}

The isotropic multiphonon rates of this work were validated against the framework of
Campbell-Deem \emph{et al.}~\cite{CampbellDeem2022} as implemented in the public package
DarkELF~\cite{Knapen2022DarkELF}, reproducing the rate spectra of
Ref.~\cite{CampbellDeem2022} (Fig.~\ref{fig:lin_comparison}).

The directional calculation was compared with Stratman and Lin~\cite{StratmanLin2024},
who computed the daily modulation of multiphonon scattering in Al$_2$O$_3$ and SiC using
the same incoherent framework with a density-of-states tensor. For this comparison our
calculation was repeated with their conventions:
\begin{itemize}
  \item halo parameters $v_0=220$, $v_E=240$, $v_{\rm esc}=500$\,km/s;
  \item no energy resolution;
  \item $c$ axis parallel to $\vec v_E$ at $t=0$ with $\theta_e = 42^\circ$, i.e.\
  $\cos\theta_E(t)=\cos^2\theta_e+\sin^2\theta_e\cos(2\pi t/24\,\mathrm{h})$;
  \item modulation amplitude $A_{\rm mod} = |R(0)-\langle R\rangle_t|/\langle R\rangle_t$.
\end{itemize}
For the heavy (contact) mediator we obtain, at $m_\chi=10$\,MeV,
$A_{\rm mod} = 2.7\%$, $5.3\%$ and $7.4\%$ for $E_{\rm th}=120$, $200$ and $300$\,meV.
At $300$\,meV the maximum is $9.8\%$ for $m_\chi\simeq 3$--$5$\,MeV, and the modulation
falls below $0.5\%$ above $30$\,MeV. Ref.~\cite{StratmanLin2024} reports modulation
amplitudes of up to $\simeq 11\%$ at $m_\chi = 10$\,MeV for $E_{\rm th}=300$\,meV,
amplitudes that grow with threshold, and a significant modulation only for
$m_\chi\lesssim 30$--$40$\,MeV. Our calculation reproduces the magnitude, the threshold
dependence and the mass range of the effect. At $m_\chi=10$\,MeV our amplitude is about
30\% smaller. This difference is within what is expected from the remaining differences
between the two calculations:
\begin{itemize}
  \item independent first-principles force constants and DOS tensors;
  \item our average over the small azimuthal dependence of the Al$_2$O$_3$ response;
  \item their evaluation of the impulse-approximation width along $\hat z$;
  \item their lower cutoff $q_{\rm BZ}=0.95$\,keV on the momentum integral.
\end{itemize}

\subsection{Limitations}

The incoherent approximation includes anisotropy only through the per-atom density of
states. It neglects the directional structure of the phonon dispersion, which dominates
single-phonon scattering at momentum transfers comparable to the inverse lattice
spacing. Dedicated single-phonon calculations in polar crystals find larger modulations
at sub-MeV masses~\cite{Griffin2018Directional}. The results presented here should
therefore be read as a conservative estimate of the low-mass modulation. The CaWO$_4$
values inherit the uncertainty of its low-frequency phonon spectrum
(Sec.~\ref{sec:tensor_validation}). Finally, the directional dependence of the
solar-neutrino background, which modulates with the solar rather than the sidereal day,
has not been included in the likelihood.

% ----------------------------------------------------------------------------
%  TODO: Fig.~\ref{fig:lin_comparison} is the DarkELF / Campbell-Deem comparison plot
%  you already made; adjust the label and the figure number of Ref. [CampbellDeem2022]
%  that you reproduced. Check the section labels used in \ref{} against your chapter.
% ----------------------------------------------------------------------------

% ---------------- BibTeX entries (move to your .bib) ------------------------
% @article{CampbellDeem2022,
%   author  = {Campbell-Deem, Brian and Knapen, Simon and Lin, Tongyan and Villarama, Ethan},
%   title   = {Dark matter direct detection from the single phonon to the nuclear recoil regime},
%   journal = {Phys. Rev. D}, volume = {106}, pages = {036019}, year = {2022},
%   eprint  = {2205.02250}, archivePrefix = {arXiv}}
%
% @article{Knapen2022DarkELF,
%   author  = {Knapen, Simon and Kozaczuk, Jonathan and Lin, Tongyan},
%   title   = {{DarkELF}: A python package for dark matter scattering in dielectric targets},
%   journal = {Phys. Rev. D}, volume = {105}, pages = {015014}, year = {2022},
%   eprint  = {2104.12786}, archivePrefix = {arXiv}}
%
% @article{StratmanLin2024,
%   author  = {Stratman, Connor and Lin, Tongyan},
%   title   = {Daily modulation of low-energy nuclear recoils from sub-GeV dark matter},
%   year    = {2024}, eprint = {2411.03433}, archivePrefix = {arXiv}}
%
% @article{Griffin2018Directional,
%   author  = {Griffin, Sin{\'e}ad and Knapen, Simon and Lin, Tongyan and Zurek, Kathryn M.},
%   title   = {Directional detection of light dark matter with polar materials},
%   journal = {Phys. Rev. D}, volume = {98}, pages = {115034}, year = {2018},
%   eprint  = {1807.10291}, archivePrefix = {arXiv}}
